%-------------------------------------------------------------------------------
\section{Background and Related Work}
\label{sec:background}
%-------------------------------------------------------------------------------

This section reviews the domain background and the prior work that motivate our design. We describe the Web PKI normative-source stack and its multi-tier constraints on certificate format, summarize the certificate-linting tools whose limitations motivate this study, and survey related work on specification extraction, knowledge-graph retrieval, and LLM output control.

\subsection{X.509 Certificate and Web PKI Normative Sources}
\label{sec:bg-pki}

X.509 certificates are the primary data representation in the Web PKI. Their structure and validation semantics are defined across layered normative sources and policy frameworks. IETF~RFC~5280 defines the foundational layer, including the certificate ASN.1 structure, field-level semantics, and the path-validation algorithm. Above it, the CA/Browser Forum Baseline Requirements (BR) impose operational and configuration constraints on commercial CAs that issue publicly trusted TLS certificates. The ETSI~EN\,319 series adds EU-specific configuration profiles aligned with the eIDAS trust framework, and the Mozilla Root Store Policy (MRSP) occupies the browser-policy layer with its own technical requirements. Together, these four normative sources form the Web PKI normative-source stack. The layers cross-reference one another and sometimes impose local overrides. Chapter~7 of CABF~BR, for example, references RFC~5280 directly while tightening individual field-level constraints.

A certificate is a signed ASN.1 object whose core payload is the tbsCertificate structure. This payload contains identity and validity metadata, including version, serialNumber, signature, issuer, validity, subject, and subjectPublicKeyInfo. Version~3 certificates further carry an extension sequence, which is where many security-relevant Web PKI constraints reside. Common examples include subjectAltName for DNS names and email names, basicConstraints for CA status and path-length limits, keyUsage and extendedKeyUsage for allowed cryptographic purposes, certificatePolicies for policy identifiers, and authority-information or revocation pointers such as AIA and CRL distribution points. These fields are encoded as DER bytes, but normative rules are usually written as natural-language statements over semantic field names, profiles, and certificate roles.

\subsection{Certificate Linting}
\label{sec:bg-linting}

Web PKI certificate issuance compliance is checked primarily with certificate-linting tools. These certificate-linting tools inspect the encoded contents of individual certificates and do not require external validation state. zlint~\cite{kumar2018tracking} is the most widely deployed certificate linter. pkilint~\cite{pkilint}, certlint~\cite{certlint}, and x509lint~\cite{x509lint} extend this ecosystem further.

zlint organizes checks as individually registered lints with metadata such as citation, effective date, and severity; each lint implements a predicate over parsed certificate structures and is commonly exercised through pass/error test certificates.

This design has two important implications for our work. First, a certificate lint is not a direct copy of a normative source sentence. It is an engineering artifact that embeds a chosen field path, a predicate, a severity level, and often a scope condition. Second, the absence of a certificate lint does not necessarily mean that no normative rule exists; it may mean that the normative rule has not been recognized as lintable or has not yet been implemented.

\subsection{Related Work}
\label{sec:bg-related}

PolicyLint~\cite{andow2019policylint} and PrivacyFlash~Pro~\cite{zimmeck2021privacyflash}
analyze privacy text~\cite{harkous2018polisis,amaral2022ai}, Hassani et al.~\cite{hassani2024rethinking} use LLMs for
legal-compliance analysis~\cite{sun2025compliance,masoudifard2024leveraging,chung2025graphcompliance}, and N-Check~\cite{feng2024analyzing} formalizes
normative rules as FOL$^*$ formulas for well-formedness analysis. These systems
extract or formalize obligations, but do not map Web PKI issuance rules to
DER-encoded certificate fields or determine their single-certificate lintability.

Specification analysis covers protocol states, formats, and code~\cite{pacheco2022automated,yen2021sage}.
PROSPER~\cite{sharma2023prosper}, ParCleanse~\cite{zheng2025validating}, and
SpecGPT~\cite{zhang2025automated} extract RFC state machines, protocol formats,
and 3GPP state-machine behavior. Wu et al.~\cite{wu2025uncovering} align RFC text
with kernel code using GraphRAG-style context~\cite{lewis2020retrieval,edge2024local,peng2025graph},
while RFCAUDIT~\cite{zheng2025rfcaudit} uses hierarchical code-semantic indexing
and retrieval-guided LLM analysis to find inconsistencies between RFC rules and
protocol implementations. SPEC2CODE~\cite{wang2025spec2code} maps fine-grained
RFC rules to functions, ProtocolGuard~\cite{song2026protocolguard} combines rule
extraction with static and dynamic verification of existing implementations, and
LLMs Unleashed~\cite{long2026llms} generates complete protocol implementations.
Unlike our pipeline, these systems do not jointly provide exhaustive rule recall,
schema-constrained IR parsing, deterministic lintability determination, and
restricted lint-predicate generation.

Output-control and verification work addresses complementary layers.
XGrammar~\cite{dong2025xgrammar} and CRANE~\cite{banerjee2025crane} constrain LLM
outputs through grammars, although constrained decoding can impair reasoning
quality~\cite{schall2025hidden,park2024grammar}. ARMOR~\cite{debnath2024armor} verifies the
RFC~5280 path-validation algorithm in Agda, while EverParse~\cite{ramananandro2019everparse}
and Nail~\cite{bangert2014nail} generate verified parsers and serializers from
declarative format grammars. They improve output form, algorithm assurance~\cite{debnath2021reengineering}, or
format handling rather than translating natural-language issuance rules into
auditable certificate lints.

Certificate studies provide the closest empirical context~\cite{durumeric2013analysis,georgiev2012most,amann2017mission,ma2021tracing,zhang2021rusted}. Kumar et
al.~\cite{kumar2018tracking} measured misissuance at ecosystem scale, and Zhang et
al.~\cite{zhang2025unicert} identified Unicode-related issuance and parsing
non-compliance, but both begin with known lints or defect classes.
Frankencerts~\cite{brubaker2014frankencerts} uses adversarial certificates for
differential testing, and SymCerts~\cite{chau2017symcerts} uses symbolic execution
to expose validator noncompliance. These systems exercise validation code rather
than derive lintable issuance rules and lint predicates from normative text.
